\chapter{Практика}


\section{Настройка интерфейсов роутера}

\begin{lstlisting}[language=C, caption={Пример настройки интерфейсов роутера}]
#R1
configure

interface gigabitethernet 1/0/1
ip address 192.168.100.1/30
ip firewall disable
exit

interface gigabitethernet 1/0/2
ip address 172.16.100.1/30
ip firewall disable
exit

1interface gigabitethernet 1/0/3
ip address 192.168.0.1/23
ip firewall disable
exit

do commit
do confirm

#R2
configure

interface gigabitethernet 1/0/1
ip address 10.10.100.2/30
ip firewall disable
exit

interface gigabitethernet 1/0/2
ip address 192.168.100.2/30
ip firewall disable
exit

1interface gigabitethernet 1/0/3
ip address 10.10.0.1/23
ip firewall disable
exit

do commit
do confirm

#R3
configure

interface gigabitethernet 1/0/1
ip address 172.16.100.2/30
ip firewall disable
exit

interface gigabitethernet 1/0/2
ip address 10.10.100.1/30
ip firewall disable
exit

interface gigabitethernet 1/0/2
ip address 192.16.0.1/23
ip firewall disable
exit

do commit
do confirm
\end{lstlisting}

\section{Настройка ip на ПК}

\begin{lstlisting}[language=C, caption={Пример настройки ip на ПК}]
#VPC1
ip 192.168.1.0/23 192.168.0.1
#VPC2
ip 10.10.1.0/23 10.10.0.1
#VPC3
ip 172.16.1.0/23 172.16.0.1
\end{lstlisting}

\section{Запуск RIP и анонс сети}

\begin{lstlisting}[language=C, caption={Пример запуска RIP и анонса сети}]
Router-left# configure

# создание процесса RIP
Router-left(config)# router rip
# анонсируем сеть (рассылаем остальным)
Router-left(config-rip)# network 192.168.0.0/23
# запускаем RIP
Router-left(config-rip)# enable

Router-left(config-rip)# exit
Router-left(config)# do commit
Router-left(config)# do confirm

Router-right# configure
Router-right(config)# router rip
Router-right(config-rip)# network 10.10.0.0/23
Router-right(config-rip)# enable

Router-right(config-rip)# exit
Router-right(config)# do commit
Router-right(config)# do confirm

Router-up# configure
Router-up(config)# router rip
Router-up(config-rip)# network 172.16.0.0/23
Router-up(config-rip)# enable

Router-up(config-rip)# exit
Router-up(config)# do commit
Router-up(config)# do confirm
\end{lstlisting}

\section{Изменения в таблице маршрутизации}

\begin{figure}[H] 
    \centering
    \includegraphics[width=0.75\textwidth]{rt.png}
    \caption{Таблица маршрутизации}
\end{figure}

У каждого маршрутизатора появится в таблице по 2 новых записи. У записей тип R, значит, они получены через RIP. \\


\section{Реакция RIP на изменение топологии}

Протоколы динамической маршрутизации обязаны реагировать на изменения в сети (добавление/удаление новых устройств, соединений).
Чтобы проверить это запустим команду ping -t, которая будет выполняться бесконечно, и удалим соединение R1-R2.
Пинг будет сыпать ошибки.
Если RIP работает корректно, то через какое-то время (по дефолту 30 секунд) сеть должна восстановиться и пинг будет проходить успешно.


\begin{figure}[H] 
    \centering
    \includegraphics[width=0.75\textwidth]{ping.png}
    \caption{Ping}
\end{figure}

Можно заметить, что пинг сначала проходил, потом стал сыпать ошибки (удалили соединение), а потом снова стал проходить.
Из этого можно сдедать вывод о том, что RIP работает корректно.





\endinput